% =============================================================================
% Project : THE_PHILOSOPHY_OF_VIABILITY
% Book    : Book I — Ontology of Viability
% Book ID : VS.B01
% Part    : Part I — Identity and Manifestation (VS.B01.P01)
% Chapter : Chapter 03 — Unity and Plurality
% Chapter ID : VS.B01.P01.C03
%
% File    : 4_ASSETS/diagrams/VS.B01.P01.C03.S02.SS10.P002.DG01.scope_of_unity_of_identity.tex
%
% Section / Chapter / Asset:
%   Diagram DG01 — Scope of Unity of Identity
%
% Structural Identifier : VS.B01.P01.C03.S02.SS10.P002.DG01
% Asset Type : DG
%
% Purpose :
%   To represent the analytical requirements for examining a claim concerning
%   unity of identity.
%
% Representational limit :
%   The subject, respect, and scope of the claim must be specified before its
%   ontological adequacy can be assessed.
%
% Dependencies :
%   TikZ; libraries used: arrows.meta, positioning, fit. All are loaded by
%   the project preamble; no per-file package or library loading.
%
% =============================================================================

\begin{figure}[htbp]
    \centering
    \hyphenpenalty=10000 \exhyphenpenalty=10000\relax  % no hyphenation inside diagram labels

    \begin{tikzpicture}[
        font=\small,
        node distance=0.4cm,
        % --- base box ---
        box/.style={
            draw, rounded corners=3pt, align=center,
            inner sep=6pt, line width=0.7pt
        },
        % --- node roles ---
        reqbox/.style={box, text width=3.0cm, minimum height=2.3cm,
                       fill=blue!8, draw=blue!55!black},
        assessbox/.style={box, text width=6.0cm, minimum height=1.3cm,
                          fill=blue!16, draw=blue!55!black},
        resultbox/.style={box, text width=7.0cm, minimum height=1.3cm,
                          fill=green!8, draw=green!40!black},
        % --- frames and labels ---
        vsgroup/.style={draw=black!45, dashed, rounded corners=5pt,
                        inner sep=8pt, line width=0.5pt},
        grouplabel/.style={font=\footnotesize\itshape, text=black!75,
                           align=left},
        % --- arrow (methodological order only) ---
        directed/.style={-{Stealth[length=2.5mm]}, line width=0.9pt}
    ]

    % ---------------------------------------------------------------
    % Three requirements (equal rank, specified together)
    % ---------------------------------------------------------------
    \node[reqbox] (subject) {
        \textbf{Subject}\\[3pt]
        What is claimed to be unified?
    };
    \node[reqbox, right=of subject] (respect) {
        \textbf{Respect}\\[3pt]
        In what ontological respect is unity attributed?
    };
    \node[reqbox, right=of respect] (scope) {
        \textbf{Scope}\\[3pt]
        Under what conditions and limits?
    };

    \node[vsgroup, fit=(subject)(respect)(scope)] (reqs) {};
    \node[grouplabel, anchor=south west] at (reqs.north west) {
        Specification of the claim
    };

    % ---------------------------------------------------------------
    % Assessment and conclusion
    % ---------------------------------------------------------------
    \node[assessbox, below=0.9cm of reqs] (assessment) {
        \textbf{Ontological Assessment}\\[3pt]
        Examine the relevant relations and supporting argument
    };
    \node[resultbox, below=0.9cm of assessment] (conclusion) {
        \textbf{Proportionate Conclusion}\\[3pt]
        A unity claim limited to its established subject,
        respect, and scope
    };

    \draw[directed] (reqs)       -- (assessment);
    \draw[directed] (assessment) -- (conclusion);

    \end{tikzpicture}

    \caption{Analytical scope of a claim concerning unity of identity. The
    subject, respect, and scope of the claim must be specified
    before its ontological adequacy can be assessed.}
    \label{fig:scope-of-unity-of-identity}
\end{figure}

% =============================================================================
% END OF DIAGRAM DG01
% =============================================================================
